\section{The static base and its conformal compactification}
\label{st-section}

We continue under the connected-horizon equality hypotheses.  Write
\[
 A_0=\Area_g(S)=4\pi r_h^2,
 \qquad m=m_{\mathrm{AH}}=\frac{r_h+r_h^3}{2},
 \qquad \kappa=\frac{1+3r_h^2}{2r_h}.
\]
The asymptotic spatial chart is balanced, so that the original metric
mass covector is $(m,0,0,0)$.  The preceding section constructed the
smooth functions and fields $u,X,N,\lambda$ and the static metric $h$ on
the original open exterior.  In particular,
\begin{equation}
 \begin{split}
 N&=u^2-|X|_g^2>0,\\
 h&=g+N^{-1}X^\flat\otimes X^\flat,
 \qquad \lambda=\sqrt N,\\
 \Delta_h\lambda&=3\lambda,
 \qquad \Hess_h\lambda=\lambda(\Ric_h+3h),
 \qquad R_h=-6.
 \end{split}
 \label{st-static-equations}
\end{equation}
Here and below the flat symbol on $X$ denotes lowering with the
original metric $g$.  The one-form $\mathcal A=X^\flat/N$ is closed,
and the open exterior is simply connected by
\cref{geo-static-quotient}.  These conclusions concern the original
exterior, rather than an auxiliary deformed manifold.

\subsection{Completion at the horizon and preservation of equality}

\begin{proposition}[The complete equality base]
\label{st-base-equality}
The metric $h$ has a smooth completion by the original boundary $S$,
with its possibly changed transverse smooth coordinate.  Its boundary
data are
\begin{equation}
 h|_{TS}=g|_{TS},\qquad
 \lambda|_S=0,\qquad
 \partial_{\nu_h}\lambda=\kappa,\qquad
 \mathrm{II}_h|_S=0.
 \label{st-base-boundary}
\end{equation}
The completed metric is complete, and $S$ is outer area-minimizing for
$h$.  Its four metric fluxes agree with those of $g$ in the original
balanced chart.  In particular,
\begin{equation}
 A_{\min,h}(S)=A_0,
 \qquad p(h)=(m,0,0,0),
 \qquad p_0(h)=\frac{r_h+r_h^3}{2}.
 \label{st-base-equality-data}
\end{equation}
For some $\tau_1>3/2$, the end satisfies
\begin{equation}
 h-b=O_{2,\alpha_1}(r^{-\tau_1}),\qquad
 \lambda-V_0=O_{2,\alpha_1}(r^{1-\tau_1}),\qquad
 h-g=O_{2,\alpha_1}(r^{-2\tau_1}).
 \label{st-base-decay}
\end{equation}
The last estimate is in physical tensor norm.
\end{proposition}

\begin{proof}
We spell out the completion because the original boundary smooth
structure is needed again in the spacetime reconstruction.  The
boundary alternatives and first jets are supplied by
\cref{eq-causal-killing,geo-boundary-collars}.

Suppose first that $u|_S>0$.  Then $N$ is an original smooth defining
function, positive into the exterior, and
\[
 X^\flat|_S=\frac1{2\kappa}\,dN|_S.
\]
In coordinates $(N,y)$ on a collar, smooth division gives
\[
 X^\flat=a(N,y)\,dN+N\beta_A(N,y)\,dy^A,
 \qquad a(0,y)=\frac1{2\kappa}.
\]
Set $N=\lambda^2$.  The coefficients of the base metric become
\begin{equation}
 \begin{split}
 h_{\lambda\lambda}&=4\lambda^2 g_{NN}+4a^2,\\
 h_{\lambda A}&=2\lambda(g_{NA}+a\beta_A),\\
 h_{AB}&=g_{AB}+\lambda^2\beta_A\beta_B,
 \end{split}
 \label{st-even-collar}
\end{equation}
where the right sides are evaluated at $(\lambda^2,y)$.  These
coefficients extend smoothly through $\lambda=0$, with even diagonal
blocks and an odd mixed block.  Reflection in $\lambda=0$ is an
isometry.  At the boundary, $h_{\lambda\lambda}=\kappa^{-2}$ and
$h_{\lambda A}=0$, proving the first three assertions in
\eqref{st-base-boundary}; its fixed hypersurface is totally geodesic.

If $u|_S=0$, let $s$ be the original inward-to-domain normal
coordinate.  The preceding boundary jets give
\[
 u=s a(s,y),\qquad X=s^2Y(s,y),\qquad a(0,y)=\kappa.
\]
Consequently, with $D=a^2-s^2|Y|_g^2$,
\[
 h=g+s^2D^{-1}Y^\flat\otimes Y^\flat,
 \qquad \lambda=s\sqrt D,
 \qquad D|_S=\kappa^2.
\]
These are smooth in the original collar.  The static Hessian equation
extends continuously to $S$ and gives $\Hess_h\lambda=0$ there.  Its
tangential part is $\kappa\mathrm{II}_h$, so the boundary is totally
geodesic.  In $h$-Gaussian distance $\rho$ it also gives
\begin{equation}
 h=d\rho^2+q(\rho),\qquad q'(0)=0,
 \qquad \lambda=\kappa\rho+O(\rho^3).
 \label{st-gaussian-boundary-jets}
\end{equation}

The completion has the same underlying topology as the original
closed exterior.  A finite-length escaping $h$-curve is a finite-length
$g$-curve because $h\geq g$.  Completeness of the original metric space
therefore gives a limit in the original closed exterior.  Interior
points are ordinary smooth base points because $N>0$ everywhere in the
interior, and points of $S$ have just been attached by regular collars.
The end has infinite $h$-length by its hyperbolic asymptotics.  Thus
there is no missing finite-length escape.

For an interior smooth enclosing cut, the quadratic-form inequality
$h\geq g$ gives its $h$-area at least its $g$-area, which is at least
$A_0$.  A cut smooth in the completed base and meeting $S$ can be
approximated, in area, by cuts whose meeting portions have been moved
a positive distance into a collar.  To see that the full-boundary
convention is preserved, apply a smooth inward collar map to the
entire exterior domain defining the cut, taper it to the identity
outside a fixed collar, and take its full intrinsic boundary.  In the
positive-lapse case perform this construction in the regular
$\lambda$ collar and keep the displacement positive before regarding
the cut as an original smooth interior cut.  Collar smoothness makes
the area error tend to zero.  Hence every completed-base cut has area
at least $A_0$, while $S$ itself has area $A_0$ and remains an
admissible full boundary.

Choose
\[
 \frac32<\tau_1<\min(2,\tau_{\mathrm{data}})
\]
slightly below the exponent in the adjoint end estimates.  Those
estimates give $u-V_0,X=O_{2,\alpha_1}(r^{1-\tau_1})$.
Since $N\sim r^2$, the identity defining $h$ gives the final estimate
in \eqref{st-base-decay}.  The identity
$\lambda-u=-|X|_g^2/(u+\lambda)$ gives the estimate for $\lambda$;
the estimate for $h-b$ follows from that for $g-b$.

Each metric-flux integrand is linear in the error and its first
background derivative and has a static-potential factor of size
$O(r)$.  The difference contributed by $h-g$ is therefore
$O(r^{3-2\tau_1})$ after integration over a coordinate sphere, and
tends to zero.  The same estimate holds for all four potentials.
The existing mass limits of $g$ consequently give the stated limits
for $h$.  This proves \eqref{st-base-equality-data} and the proposition.
\end{proof}

\subsection{The Riemannian Einstein metric}

Use a circle coordinate $\theta$ of any fixed positive period on the
end and put
\begin{equation}
 G=h+\lambda^2d\theta^2.
 \label{st-wick-metric}
\end{equation}
The warped-product Christoffel symbols involving the circle are
$\Gamma^\theta_{i\theta}=\lambda_i/\lambda$ and
$\Gamma^i_{\theta\theta}=-\lambda\nabla_h^i\lambda$.  Thus
\begin{equation}
 \Ric_G|_{T\Omega}=\Ric_h-\lambda^{-1}\Hess_h\lambda=-3h,
 \qquad
 (\Ric_G)_{\theta\theta}=-\lambda\Delta_h\lambda=-3\lambda^2,
 \qquad (\Ric_G)_{i\theta}=0.
 \label{st-wick-einstein}
\end{equation}
In particular $G$ is a four-dimensional Riemannian Einstein metric.
It is invariant under circle translation and time reflection
$\theta\mapsto-\theta$.

Only an end collar is needed for the conformal-boundary regularity
theorem below.  The following also explains why the horizon causes no
singularity if a global Riemannian filling is used.

\begin{lemma}[The regular circle bolt]
\label{st-bolt}
If the period of $\theta$ is $2\pi/\kappa$, collapsing its circles
over $S$ extends \eqref{st-wick-metric} across a regular bolt.  In
polar Cartesian coordinates it is initially a $C^{2,\alpha}$
Einstein metric for every $0<\alpha<1$ and hence admits smooth
Einstein coordinates across the bolt.
\end{lemma}

\begin{proof}
Use \eqref{st-gaussian-boundary-jets} and set
$\vartheta=\kappa\theta$, of period $2\pi$.  On each boundary
coordinate patch, write $x=\rho\cos\vartheta$ and
$y=\rho\sin\vartheta$.  The normal part is
\[
 d\rho^2+(\lambda/\kappa)^2d\vartheta^2
 =dx^2+dy^2+
 \frac{(\lambda/(\kappa\rho))^2-1}{\rho^2}
 (x\,dy-y\,dx)^2.
\]
The coefficient in front of the final squared one-form is a smooth
one-sided function of $\rho$ and the bolt variables, equal to a smooth
bolt function plus $O(\rho)$.  The squared one-form has quadratic
Cartesian coefficients.  Its $O(\rho)$ correction is consequently
$C^{2,\alpha}$ for every $\alpha<1$.  Likewise
$q(\rho)=q(0)+\rho^2q_2+O(\rho^3)$ gives a $C^{2,\alpha}$
tangential block.  There are no mixed $d\rho$--bolt terms in Gaussian
coordinates.  Positivity follows at $\rho=0$ from $q(0)>0$.

The Einstein equation holds away from the bolt and, because the
metric is $C^2$, also holds at the bolt by continuity of Ricci
curvature.  Interior Einstein regularity in harmonic coordinates
applies, for instance \cite[Theorem~5.2]{deturckkazdan1981}.
This gives smooth, indeed analytic, Einstein coordinates there.
The circle action and its fixed set are then smooth: isometries
preserve harmonic coordinates and the local Killing field satisfies
the corresponding elliptic equation.  Normal exponential coordinates
to the fixed set recover the usual smooth disk action.  Thus this is
a regular bolt, with no conical angle defect.
\end{proof}

\subsection{A weighted gauge on the hyperbolic tube}

The background for the end calculation is the complete hyperbolic
tube
\begin{equation}
 G_0=b+V_0^2d\theta^2
 =d\eta^2+\sinh^2\eta\,\sigma+\cosh^2\eta\,d\theta^2
 \quad\hbox{on }B^3\times S^1.
 \label{st-hyperbolic-tube}
\end{equation}
Its sectional curvature is $-1$.  Near infinity the defining
function $z_0=2e^{-\eta}$ gives
\begin{equation}
 G_0=z_0^{-2}\left[
 dz_0^2+(1-z_0^2/4)^2\sigma+(1+z_0^2/4)^2d\theta^2
 \right].
 \label{st-tube-compactification}
\end{equation}
We extend $z_0$ to a positive smooth function in the interior when
defining global weighted spaces.  All such extensions give equivalent
norms.  These are the intrinsic weighted H\"older norms measured in
bounded physical coordinate patches; a tensor in $C^{k,\alpha}_\delta$
has physical size and derivatives $O(z_0^\delta)$.

\begin{lemma}[Gauge and improved decay]
\label{st-weighted-gauge}
Let an Einstein metric $G$ on a tube end satisfy
$G-G_0\in C^{2,\alpha}_{\beta}$ for some $\beta>3/2$.
After a boundary-identity change of coordinates, its difference from
$G_0$ belongs to $C^{k,\alpha_0}_{\delta}$ on a smaller end for
every integer $k$ and every $0<\delta<3$.  Here one can choose
$\alpha_0>0$ below the original exponent.  The coordinate change
preserves circle translation and time reflection if $G$ does.
In particular $G$ has an ordinary $C^{2,\alpha_2}$ conformal
compactification for some $\alpha_2>0$.
\end{lemma}

\begin{proof}
Fix $\gamma$ with $3/2<\gamma<\min(2,\beta)$ and lower the
H\"older exponent slightly.  First cut $G-G_0$ off toward the interior,
leaving it unchanged on a sufficiently distant end.  Call the
resulting global metric $G_E=G_0+E$.  Because $\gamma<\beta$, the
global $C^{2,\alpha_0}_\gamma$ norm of $E$ tends to zero as the
cutoff is moved out.  The cutoff is a function of $\eta$, so retains
the stated symmetries.

Seek a harmonic map
\[
 \Phi:(B^3\times S^1,G_E)\longrightarrow(B^3\times S^1,G_0),
 \qquad \Phi(x)=\exp^{G_0}_x W(x),
\]
with $W\in C^{3,\alpha_0}_\gamma$.  Parallel translation along
the exponential segment identifies its tension field with a section
of the fixed background tangent bundle.  The negative linearization
in $W$ at $(E,W)=(0,0)$ is
\begin{equation}
 J=\nabla^*\nabla-\Ric_{G_0}=\nabla^*\nabla+3.
 \label{st-map-jacobi}
\end{equation}
It has no $L^2$ kernel, since its quadratic form is
$\int(|\nabla W|^2+3|W|^2)$.  Theorem~C(c) and Proposition~E
of \cite{lee2006}, applied to one-forms after lowering an index,
give an isomorphism in physical weights $-1<\delta<4$:
the indicial radius is
$\sqrt{3^2/4+1+3}=5/2$, and the Fredholm index is zero with
kernel equal to the $L^2$ kernel.  In particular,
\[
 J:C^{3,\alpha_0}_\gamma\longrightarrow C^{1,\alpha_0}_\gamma
\]
is an isomorphism.

The tension expression is a smooth map between these Banach spaces.
In physical unit patches its coefficients are smooth functions of
$E,G_E^{-1}=(G_0+E)^{-1}$, one derivative of $E$, and the first two derivatives of
$W$; target curvature and its derivatives are uniformly bounded.
Products of weighted errors have at least their summed weight.
Consequently its nonlinear remainder satisfies, on a fixed small
ball,
\[
 \|\mathcal R(E,W)-\mathcal R(E,\widetilde W)\|_{C^{1,\alpha_0}_\gamma}
 \leq C\bigl(\|E\|_{C^{2,\alpha_0}_\gamma}
       +\|W\|_{C^{3,\alpha_0}_\gamma}
       +\|\widetilde W\|_{C^{3,\alpha_0}_\gamma}\bigr)
       \|W-\widetilde W\|_{C^{3,\alpha_0}_\gamma}.
\]
The inverse of $J$ therefore gives a contraction, or equivalently
the implicit function theorem gives a unique small $W$, with
$\|W\|_{C^{3,\alpha_0}_\gamma}\leq C\|E\|_{C^{2,\alpha_0}_\gamma}$.
The map is a local diffeomorphism by smallness of its first jet.  Its
displacement is bounded, so it is proper for the complete background
metric.  It is homotopic to the identity through its exponential
displacements.  A proper local diffeomorphism is a covering map, and
this homotopy makes its induced fundamental-group map surjective;
the covering therefore has one sheet.  Thus $\Phi$ is a global
diffeomorphism.  In compact coordinates its displacement is
$O(z_0^{1+\gamma})$, so it restricts to the identity at infinity.
Uniqueness shows that it commutes with circle translations and time
reflection.

Set $\widehat G=(\Phi^{-1})^*G_E$ and $k=\widehat G-G_0$.
The identity map from $\widehat G$ to $G_0$ is harmonic, and
$k\in C^{2,\alpha_0}_\gamma$.  On a smaller end $\widehat G$
is Einstein.  The Ricci equation with this harmonic-map gauge has
linearization
\begin{equation}
 P=\frac12(\Delta_L+6)
   =\frac12\bigl(\nabla^*\nabla-2\mathring R_{G_0}\bigr)
 \label{st-einstein-linearization}
\end{equation}
at $G_0$; see \cite[Equations~(1.2)--(1.3)]{lee2006}.
In each fixed-size physical patch the gauged Ricci expression is
uniformly elliptic.  Differentiating its difference from the
background equation and applying interior Schauder estimates on
nested patches gives successively
\[
 \|k\|_{C^{j+2,\alpha_0}(B_1)}
 \leq C_j\bigl(\|k\|_{C^{j+1,\alpha_0}(B_2)}
                 +\|k\|_{C^{2,\alpha_0}(B_2)}\bigr).
\]
The constants are uniform along the end because the background has
bounded geometry and $k$ is uniformly small in $C^{2,\alpha_0}$.
The defining function is comparable on each such patch.  Induction
therefore preserves the factor $z_0^\gamma$ and gives
$k\in C^{j,\alpha_0}_\gamma$ for every $j$.

Taylor expansion of the gauged equation now has the exact form
\begin{equation}
 Pk=\mathcal Q(k,\nabla k,\nabla^2k),\qquad
 \mathcal Q\in C^{j,\alpha_0}_{2\gamma}\quad\hbox{for every }j.
 \label{st-quadratic-einstein-error}
\end{equation}
Indeed its constant term vanishes at $G_0$, its first-order term is
$P$, and every remaining term contains at least two error factors,
counting derivatives as factors in physical norms.

We verify injectivity for the next use of the weighted theorem.
On the tube, $Z=\partial_\eta$ has
\[
 \operatorname{div}_{G_0}Z=2\coth\eta+\tanh\eta\geq3.
\]
For a compactly supported scalar $f$, integration by parts and
Cauchy--Schwarz give
\[
 3\int f^2\leq-2\int f\,Zf
 \leq2\|f\|_2\,\|df\|_2,
 \qquad \int|df|^2\geq\frac94\int f^2.
\]
The integration can first omit $\eta<\varepsilon$; its additional
boundary term tends to zero because the tubular area is
$O(\varepsilon^2)$.  Kato's Inequality gives the same lower bound
for $\int|\nabla T|^2$ in terms of $\int|T|^2$.
For a tracefree symmetric two-tensor on curvature $-1$,
$\mathring R T=T$, so the operator inside the parentheses in
\eqref{st-einstein-linearization} has quadratic form at least
$\frac14\|T\|_2^2$.  On a pure trace tensor,
$\mathring R(fG_0)=-3fG_0$, so its zeroth-order term is favorable.
The trace and tracefree bundles are parallel.  Thus $P$ has no
$L^2$ kernel.

Theorem~C(c) and Proposition~D of \cite{lee2006} apply to this smooth,
complete conformally compact background.  For $\Delta_L+6$ in
dimension four the indicial radius is $3/2$, so $P$ is an
isomorphism at every physical weight $0<\delta<3$.
Cut $k$ off once more toward the interior.  Its equation on the
whole tube is \eqref{st-quadratic-einstein-error} plus a smooth
compactly supported commutator.  Since $2\gamma>3$, that source
belongs to $C^{j,\alpha_0}_\delta$ for every $\delta<3$.
Solve it at any $\delta$ with $2<\delta<3$ and compare with the
cut-off $k$ at weight $\gamma$.  Their difference is in the kernel
at the smaller weight, hence vanishes.  Higher weighted regularity
follows by the same equation.  Smaller positive weights follow by
inclusion.

Finally, for an ordinary compact coordinate component of
$z_0^2k$, two coordinate derivatives have size
$O(z_0^{\delta-2})$.  The physical H\"older estimate, compared on
patches of compact diameter comparable to $z_0$, gives an ordinary
$C^{2,\alpha_2}$ extension whenever
$0<\alpha_2<\min(\alpha_0,\delta-2)$.  Equivalently this is the
weighted-to-ordinary inclusion of \cite[Lemma~3.7(b)]{lee2006},
with the tensor weight included.  This proves the claimed
preliminary conformal compactness.
\end{proof}

\subsection{Smooth infinity and the normalized end}

\begin{proposition}[The normalized smooth expansion]
\label{st-end-regularity}
There is a geodesic defining function $z$ for the conformal
representative $d\theta^2+\sigma$ and a boundary-identity spatial
collar identification in which
\begin{equation}
 \begin{split}
 h&=z^{-2}(dz^2+H(z)),\\
 H(z)&=\sigma-\frac12z^2\sigma+z^3H_3+O(z^4),\\
 z\lambda&=1+\frac14z^2+z^3\lambda_3+O(z^4).
 \end{split}
 \label{st-normalized-expansion}
\end{equation}
All functions and tensor coefficients are smooth up to $z=0$.
The full third coefficient of the compactified four-metric is
\begin{equation}
 G_3=H_3+2\lambda_3d\theta^2,
 \qquad \operatorname{tr}_{d\theta^2+\sigma}G_3=0,
 \qquad \operatorname{div}_{d\theta^2+\sigma}G_3=0.
 \label{st-free-tt-data}
\end{equation}
The geodesic gauge preserves the circle factor and time reflection.
The original and new defining radii are comparable, their collar
maps extend to the same conformal boundary, and the resulting
spatial mass comparison is covered by
\cref{lem:setup-mass-covariance}.
\end{proposition}

\begin{proof}
Equations \eqref{st-base-decay} and \eqref{st-wick-metric} give
$G-G_0\in C^{2,\alpha_1}_{\tau_1}$ in physical four-dimensional
tensor norm: the relative circle coefficient is
$(\lambda^2-V_0^2)/V_0^2=O(r^{-\tau_1})$.
Apply \cref{st-weighted-gauge}.  The resulting metric is smooth
in the interior, Einstein with $\Ric=-3G$, has an ordinary $C^2$
conformal compactification, and has the smooth conformal
representative $d\theta^2+\sigma$.

These are the hypotheses of
\cite[Theorem~A]{cdls2005}.  That theorem is local to a collar of a
compact conformal boundary and does not require a global bolt
filling.  It supplies a boundary-identity collar diffeomorphism
placing the metric in geodesic form.  Its compactification is smooth
in even total dimension, hence in the present dimension four.
The defining function is chosen for the specified representative.
Both its defining eikonal equation and the inward normal flow
are invariant under circle translation and reflection; uniqueness
with that boundary normalization makes the collar map equivariant.
Equivariance under translations permits at most a spatially
dependent circle shift, and reflection together with boundary
identity forces that shift to vanish.  Thus the geodesic gauge is a
spatial one and has the block form in
\eqref{st-normalized-expansion}.

For completeness we compute the coefficients needed later.  Write
\[
 G=z^{-2}(dz^2+\gamma(z)),\qquad
 \gamma(0)=\gamma_0=d\theta^2+\sigma.
\]
The tangential Einstein equation and its mixed constraint are
\begin{align}
 0={}&z\gamma''-z\gamma'\gamma^{-1}\gamma'
  +\frac z2\operatorname{tr}(\gamma^{-1}\gamma')\gamma'
  -2\gamma'-\operatorname{tr}(\gamma^{-1}\gamma')\gamma
  -2z\Ric_{\gamma},
 \label{st-geodesic-einstein-equation}\\
 0={}&\operatorname{div}_{\gamma}
       (\gamma^{-1}\gamma')
       -d\operatorname{tr}(\gamma^{-1}\gamma').
 \label{st-geodesic-mixed-constraint}
\end{align}
They follow by using the second form $\gamma'/2$ of the compact
metric $dz^2+\gamma$ and the conformal Ricci formula.  Expanding the
first equation at order zero gives $\gamma'(0)=0$.
Writing $\gamma=\gamma_0+z^2\gamma_2+z^3\gamma_3+\cdots$, its
order-$z$ equation gives
\[
 \gamma_2+\operatorname{tr}_{\gamma_0}(\gamma_2)\gamma_0
 =-\Ric_{\gamma_0}.
\]
Here $\Ric_{\gamma_0}=\sigma$ and $R_{\gamma_0}=2$, so
\[
 \gamma_2=\frac12d\theta^2-\frac12\sigma.
\]
The order-$z^2$ equation gives
$\operatorname{tr}_{\gamma_0}\gamma_3=0$, and the same order in
\eqref{st-geodesic-mixed-constraint} gives
$\operatorname{div}_{\gamma_0}\gamma_3=0$.
Since the circle coefficient of $\gamma$ is $(z\lambda)^2$,
these statements are exactly
\eqref{st-normalized-expansion}--\eqref{st-free-tt-data}.

We verify the spatial mass comparison from the two asymptotic
representations of the same metric $h$. In the original spatial chart,
\eqref{st-base-decay} gives $h-b=O_2(r_{\mathrm o}^{-\tau_1})$
with $\tau_1>3/2$. In the new geodesic chart put $r=z^{-1}-z/4$.
The exact identities
\[
 \frac{dr^2}{1+r^2}=z^{-2}dz^2,\qquad
 r^2\sigma=z^{-2}(1-z^2/4)^2\sigma
\]
and the smooth expansion \eqref{st-normalized-expansion} give
$h-b=O_j(r^{-3})$ for every fixed differentiated order $j$.
Also $R_h=-6$ by the static equations. In each chart, therefore,
the weighted quadratic integrability of
\cite[Equation~(2.8a)]{chruscielherzlich2003} follows from
$\int^\infty r^{2-2q}\,dr<\infty$ with $q>3/2$;
its scalar-curvature condition (2.8b) is automatic.
The frame decay is $o(r^{-3/2})$, as required by (2.13), and
uniform metric equivalence holds sufficiently far out.
Theorem~2.3 of that paper identifies the metric mass functionals
up to a hyperbolic background isometry. The transition estimates
are conclusions of that theorem, so differentiated estimates for
the geodesic coordinate map are not an additional input here.
Both the harmonic-map identification and the geodesic collar
identification induce the identity on the spatial conformal boundary.
The resulting hyperbolic isometry acts trivially on $\mathbb S^2$
at infinity and is therefore the identity. This proves the asserted
mass comparison.

For radial comparability, write $z_{\mathrm o}$ for the original
tube defining function. The original normalization gives
$z_{\mathrm o}\lambda\to1$, while
\eqref{st-normalized-expansion} gives $z\lambda\to1$.
The collar changes preserve the same unrescaled lapse, hence
$z/z_{\mathrm o}\to1$ at corresponding points. The harmonic-map
displacement gives $z_0/z_{\mathrm o}\to1$ in the intermediate
gauge, so also $z/z_0\to1$.
\end{proof}

The expansion fixes the decay and normalization needed to evaluate
the static Heintze--Karcher deficit in the next section. Its third-order
trace records the mass, while the horizon area and surface gravity
supply the inner boundary term.
